\chapter{Spring Kafka: Producers, Listeners, Poison Pills}
\label{ch:springkafka}

\section{Producing}

Publishing an event in this project is two moving parts: a
\code{KafkaTemplate} and configuration choosing serializers. From
auth-service:

\begin{lstlisting}[language=Java]
kafkaTemplate.send("ts.user.registered", event.userId(), event);
//                  topic                key             value
\end{lstlisting}

\begin{lstlisting}[language=yaml]
spring:
  kafka:
    bootstrap-servers: localhost:9092
    producer:
      key-serializer: org.apache.kafka.common.serialization.StringSerializer
      value-serializer: org.springframework.kafka.support.serializer.JsonSerializer
\end{lstlisting}

The key (\code{userId}) is a string; the value is the event record
serialized to JSON. Keying by user id is the ordering decision from
Chapter~10 made concrete. \code{send()} is asynchronous --- it returns a
future; the message sits in a client-side buffer until the broker acks.
Two operational consequences: a service can \emph{start} without a broker
(sends fail later, per call --- which is why auth-service could deploy to
the cluster before Kafka), and a service that crashes right after
\code{send()} returns may lose the buffered message --- the
outbox pattern below exists for when that matters.

\section{Consuming}

\begin{lstlisting}[language=Java]
@KafkaListener(topics = "ts.user.registered",
               groupId = "notification-group")
public void on(UserRegisteredEvent event) {
    mailService.sendActivation(event);
}
\end{lstlisting}

Spring runs a polling loop per listener container, deserializes each
record, invokes the method, and commits offsets after successful
processing --- at-least-once out of the box. The consumer half of the
YAML holds this chapter's most production-relevant lines:

\begin{lstlisting}[language=yaml]
spring:
  kafka:
    consumer:
      group-id: course-group
      auto-offset-reset: earliest                    (*@\co{1}@*)
      key-deserializer: ...ErrorHandlingDeserializer (*@\co{2}@*)
      value-deserializer: ...ErrorHandlingDeserializer
      properties:
        spring.deserializer.value.delegate.class: ...JsonDeserializer
        spring.json.trusted.packages: com.ts.common.dto  (*@\co{3}@*)
\end{lstlisting}

\begin{description}
  \item[\co{1}] Where a \emph{brand-new} group starts when it has no
  committed offset: \code{earliest} = from the beginning of retention;
  the default \code{latest} = only new messages. \code{earliest} means a
  consumer deployed late still processes events published before it
  existed --- the right choice when events are facts to act on.
  Irrelevant once the group has offsets.
  \item[\co{2}] The poison-pill defense --- next section.
  \item[\co{3}] JSON deserialization instantiates classes named in
  message headers; trusting only \code{com.ts.common.dto} stops a
  malicious or buggy producer from making the consumer instantiate
  arbitrary classes. A security boundary, not boilerplate.
\end{description}

\section{The poison pill}

A \emph{poison pill} is a message that fails deserialization forever ---
a producer bug, a renamed field, a corrupted payload. Without defenses
the listener throws \emph{before} the message can be skipped, the offset
never advances, and the container retries the same record endlessly:
the partition is stalled and the logs fill at full speed.

\code{ErrorHandlingDeserializer} wraps the real deserializer: on
failure it emits a typed error instead of throwing inside the poll loop,
letting Spring's error handler skip past the record (optionally after
retries, optionally publishing it to a dead-letter topic). The
composition --- error-handler outside, JSON delegate inside --- is the
standard production pattern, and it is why the config names \emph{two}
deserializer classes per side.

\begin{pitfall}[one bad message, total silence]
Symptom: notification mails stop entirely; the consumer logs the same
deserialization stack trace thousands of times. Cause: a poison pill on
one partition with no error-handling deserializer --- the loop can
neither process nor pass it. This failure hides in systems that never
tested a malformed message. Verify your defense once, deliberately:
publish garbage to a dev topic and watch it get skipped.
\end{pitfall}

\section{Advanced corners}

\begin{description}
  \item[Dead-letter topics.] \code{DefaultErrorHandler} +
  \code{DeadLetterPublishingRecordRecoverer}: after N failed attempts,
  the record lands on \code{<topic>.DLT} for inspection and replay ---
  failure handling with an audit trail.
  \item[The outbox pattern.] ``Save the user \emph{and} publish the
  event'' is two systems and cannot be atomic directly. The outbox:
  write the event into an \code{outbox} table in the \emph{same database
  transaction} as the user row; a relay (or Debezium reading the WAL)
  publishes from that table. If auth-service ever must not lose a
  registration event, this is the answer.
  \item[Schema registry.] Chapter~6's shared event classes solve
  producer/consumer agreement inside one repo. At organization scale the
  contract moves into a schema registry (Avro/Protobuf schemas, checked
  for compatibility at publish time) --- the same idea, enforced by
  infrastructure instead of a shared jar.
\end{description}

\begin{quiz}
\begin{enumerate}
  \item Why can auth-service start when Kafka is down, and at what
  moment does a Kafka outage become visible to it?
  \item A new consumer group is deployed today. Explain what
  \code{auto-offset-reset: earliest} makes it do, and when this setting
  stops having any effect for that group.
  \item Walk through the poison-pill failure without
  \code{ErrorHandlingDeserializer}, then with it. Where does the bad
  record end up in a full production setup?
  \item Why is \code{spring.json.trusted.packages} a security control?
  \item When is the outbox pattern necessary, and what pair of writes
  does it make effectively atomic?
\end{enumerate}
\end{quiz}
